\section{The background and the Reiterer--Trubowitz formulation}\label{sec:setting}

This section fixes notation. \Cref{sec:setting-background,sec:setting-system} recall the parts
of~\cite{RT2019} that we use, with their notation wherever possible. \Cref{sec:setting-linear,sec:setting-sectors,sec:setting-constraints,sec:setting-realizations}
define the objects of the Main Theorem: the linearized family $L(s)$, the constraint map $C(s)$, the
constraint propagation operator $K(s)$, the Floquet classes and the function spaces.
Section numbers in references to~\cite{RT2019} are those of its arXiv version 1203.3766v1, whose source we
used throughout.

\subsection{The Choptuik spacetime}\label{sec:setting-background}

Let
\[
  M=\{(u_-,u_+)\in\R^2:\ u_+<0,\ u_+<u_-<-u_+/81\},\qquad
  \xi=\frac{u_+-u_-}{u_++u_-}.
\]
Reiterer and Trubowitz~\cite[Main result]{RT2019} prove the existence of constants $K,\mu>0$, known to
$80$ digits ($K\approx1.72272620$, $\mu\approx0.16830708$), and of real analytic functions
$\phi,\zeta,Q$ on $M$ with $\mu+Q\xi^2>0$, such that the spherically symmetric metric
\begin{equation}\label{eq:metric}
  \gbar=e^{-2\zeta}\Big(-\frac{2\,(du_-\otimes du_++du_+\otimes du_-)}{\mu^2(u_++u_-)^2}
  +\frac{\xi^2}{(\mu+Q\xi^2)^2}\,g_{S^2}\Big)
\end{equation}
on $M\times S^2$ and the scalar field $\phibar=\phi$ solve
$\operatorname{Ric}_g=2\,d\phi\otimes d\phi$, $\Box_g\phi=0$. The map
$\Theta(u_-,u_+)=e^{-2\pi\mu}(u_-,u_+)$ satisfies
\[
  \Theta^*\gbar=e^{-2K}\gbar,\qquad \phibar\circ\Theta=-\phibar,
\]
and the scalar curvature is not identically zero, so the spacetime is not flat and blows up at the
singular point $(u_-,u_+)=(0,0)$, which is not in $M$. The boundary $u_+=u_-$ is the axis of symmetry, a
removable coordinate singularity. The null hypersurface $u_-=0$ is the past light cone of the singular
point; $M$ contains a neighbourhood of it.

\emph{Coordinates.} We use the coordinates $(\tau,\xi)$ of~\cite{RT2019}, defined by
\begin{equation}\label{eq:tauxi}
  u_\sigma=-(1+\sigma\xi)\,e^{-\mu\tau},\qquad\sigma=\pm.
\end{equation}
In them $\Theta$ is the translation $\tau\mapsto\tau+2\pi$, the axis is $\xi=0$, the light cone $u_-=0$ is
$\xi=1$, and the condition $u_-<-u_+/81$ becomes $\xi<41/40$. Replacing $\xi\ge0$ by the Cartesian
variable $x\in\R^3$ with $|x|=\xi$, the spacetime becomes the open set
\[
  \Mhat=\R_\tau\times B_{41/40}(0)\subset\R\times\R^3,
\]
on which $\gbar$ and $\phibar$ are real analytic, including at the axis $x=0$ and across the light cone
$|x|=1$~\cite{RT2019}. All perturbations below are defined on $\Mhat$.

\emph{Logarithmic time.} Since $\Theta$ rescales lengths by $e^{-K}$ and acts by $\tau\mapsto\tau+2\pi$,
the time $T=K\tau/2\pi$ is the usual logarithmic time of critical collapse, in which $\Theta$ is a
translation by $K$. The scalar field returns to itself under $\Theta^2$, so the echoing period is
$\Delta=2K=3.44545240\ldots$, in agreement with the numerical value $3.445452402(3)$ of
Mart\'in-Garc\'ia and Gundlach~\cite{MartinGarciaGundlach2003}, as already noted in~\cite{RT2019}.

\subsection{Frame variables and the system}\label{sec:setting-system}

Reiterer and Trubowitz write the equations as a first order system, quadratically nonlinear, for the
constant $\mu$ and four functions $\omega=(\omega_1,\omega_2,\omega_3,\omega_4)$ of $(\tau,\xi)$. With the
radial null vector fields
\[
  D^\sigma=\sigma\,\partial_\tau+\mu(1+\sigma\xi)\,\partial_\xi,\qquad\sigma=\pm,
\]
and the reflection $(Pf)(\tau,\xi)=f(\tau,-\xi)$, they are
\begin{equation}\label{eq:omega}
\begin{aligned}
  \omega_1&=2\xi Q+(D^-+D^+)\big(\zeta+\log(\mu+\xi^2Q)\big),&
  \omega_2^\sigma&=D^\sigma\big(\zeta+\log(\mu+\xi^2Q)\big)-\sigma\mu,\\
  \omega_3^\sigma&=D^\sigma\zeta-\sigma\mu,&
  \omega_4^\sigma&=D^\sigma\phi,
\end{aligned}
\end{equation}
with $\omega_i^-=\omega_i$ and $\omega_i^+=-P\omega_i$ for $i=2,3,4$~\cite[\S2]{RT2019}. The
Einstein--scalar field equations are equivalent to the vanishing of ten explicit functions
$\Omega_i^\sigma(\mu,\omega)$, $i\in\{1,2,2\sharp,3,4\}$, each of the form
\[
  \Omega_i^\sigma=\xi\,(D^\sigma+\sigma\mu)(\text{one component of }\omega)
  +\mu\,(\text{linear in }\omega)+\xi\,(\text{quadratic in }\omega).
\]
The reflection symmetry reduces them to the five equations $\Omega^+=0$, posed on the cylinder
$(\R/4\pi\Z)\times(-\xi_*,\xi_*)$, $\xi_*>1$, for functions with the symmetries
\begin{equation}\label{eq:symmetries}
  \omega_1,\omega_2,\omega_3\ \text{$2\pi$-periodic},\qquad \omega_4\ \text{$2\pi$-antiperiodic},\qquad
  P\omega_1=-\omega_1.
\end{equation}

\emph{Series.} Each $\omega_i$ is expanded as a Fourier series in $\tau$ with period $4\pi$ and a
Chebyshev series in $\xi$:
\begin{equation}\label{eq:series}
  v(\tau,\xi)=\sum_{m,n\in\Z}v_{mn}\,e^{im\tau/2}e^{in\theta},\qquad \xi=\cos\theta,\qquad
  v_{m,-n}=v_{mn}.
\end{equation}
The symmetries~\eqref{eq:symmetries} say that $\omega_1,\omega_2,\omega_3$ contain only even $m$,
$\omega_4$ only odd $m$, and $\omega_1$ only odd $n$. In these coordinates $\partial_\tau$ acts by
$im/2$, the reflection by $(-1)^n$, multiplication by $\xi$ by averaging the neighbours $n\pm1$, and
products of functions by convolution of coefficients.

\emph{The selected system and the constraints.} Reiterer and Trubowitz split the five equations
$\Omega^+=0$ into a part with as many equations as unknowns and a remainder, the constraints, by two
selection operators $S$ and $S^\sharp$: $\Omega^+=0$ if and only if $S\Omega^+=0$ and $S^\sharp\Omega^+=0$.
The selected part is
\begin{equation}\label{eq:F}
  F(\mu,\omega):=S\,\Omega^+(\mu,\omega)
  =O_\mu\,\omega+\mu\,S\Gamma_1\omega+S\Xi\,\Gamma_2(\omega,\omega),
\end{equation}
where:
\begin{itemize}[leftmargin=*]
\item $O_\mu=\partial_\tau+\mu D_O$, with $D_O=\operatorname{diag}(\xi\partial_\xi+1,\
  (1+\xi)\partial_\xi+1,\ (1+\xi)\partial_\xi+1,\ (1+\xi)\partial_\xi+1)$, is the free part;
\item $\Gamma_1$ is linear and $\Gamma_2$ is symmetric bilinear; both are built from reflections and
  pointwise products, without derivatives;
\item $\Xi$ is multiplication by $\xi$, and $S$ contains the regularized division by $\xi$,
  $R_\xi f=(f-f|_{\xi=0})/\xi$.
\end{itemize}
The explicit formulas are in~\cite[\S3]{RT2019}. We write $F^\sharp(\mu,\omega):=S^\sharp\Omega^+(\mu,\omega)$
for the constraints.

\emph{The background.} Reiterer and Trubowitz construct a solution
$(\mu,\omegabar)$ of $F=0$, $F^\sharp=0$, as a fixed point of a contraction in a weighted $\ell^1$ space
of coefficients with weights $(65/64)^{|m|}(5/4)^{|n|}$, near explicit numerical data, and from it the
spacetime of \cref{sec:setting-background}~\cite[\S\S4--5]{RT2019}. The solution is normalized by the
condition $\operatorname{Re}(\omegabar_4)_{10}=0$, which fixes the translation invariance in $\tau$. Their final
bounds are
\begin{equation}\label{eq:RTball}
  |\mu-(\mu_A+\mu_B)|\le2^{-277},\qquad \|\omegabar-(\omega_A+\omega_B)\|\le2^{-277},\qquad
  \|\omega_B\|\le2^{-25},
\end{equation}
where $(\mu_A,\omega_A)$ and $(\mu_B,\omega_B)$ are explicit data published with~\cite{RT2019}, and the
norm is their weighted $\ell^1$ norm. Most of our bounds use only $(\mu_A,\omega_A)$ and the consequence
$\|\omegabar-\omega_A\|\le2^{-25}+2^{-277}$.

\subsection{The linearized family}\label{sec:setting-linear}

A perturbation of Floquet type is $\omega=\omegabar+\epsilon\,e^{s\tau}h$, where $h$ is a $4\pi$-periodic
multifield. Here $\mu$ is held fixed, and so are the coordinates $(\tau,\xi)$: perturbations of $\mu$ would
move the null coordinates~\eqref{eq:tauxi}. That this is no loss of generality for physical modes is part
of \cref{sec:physical}. The normalization of the phase is not imposed on $h$.

\begin{definition}[The linearized family]\label{def:L}
For $s\in\C$ and a $4\pi$-periodic multifield $h$ with $Ph_1=-h_1$, set
\[
  L(s)h:=e^{-s\tau}\,D_\omega F(\mu,\omegabar)\big[e^{s\tau}h\big],\qquad
  C(s)h:=e^{-s\tau}\,D_\omega F^\sharp(\mu,\omegabar)\big[e^{s\tau}h\big].
\]
\end{definition}

\begin{lemma}[$L$ and $C$ are affine in $s$]\label{lem:affine}
$L(s)=L(0)+s\,I$ and $C(s)=C_0+s\,C_1$, where
\[
  L(0)=J+B,\qquad J=O_\mu,\qquad B=\mu\,S\Gamma_1+2\,S\Xi\,\Gamma_2(\omegabar,\cdot\,),
\]
and $C_0,C_1$ are given explicitly in \cref{app:formulas}.
\end{lemma}

\begin{proof}
In~\eqref{eq:F}, the only derivative in $\tau$ is the one in $O_\mu$, with coefficient one. The operators
$\Gamma_1$, $\Gamma_2(\omegabar,\cdot)$, $\Xi$ and $S$ act in $\xi$ and by pointwise products, so they
commute with multiplication by $e^{s\tau}$, and $e^{-s\tau}\partial_\tau e^{s\tau}=\partial_\tau+s$. In
$\Omega^+=\Xi H_\mu\omega+\mu\Gamma_1\omega+\Xi\Gamma_2(\omega,\omega)$ the derivative $\partial_\tau$ occurs
only in $H_\mu=H_1\partial_\tau+(\text{terms in }\xi)$, where $H_1$ is a constant matrix of identities and
reflections~\cite[\S\S2--3]{RT2019}. Hence the same computation gives $C(s)=C_0+sC_1$ with
$C_1=S^\sharp\Xi H_1$.
\end{proof}

A root of $L$ is a point $s_0$ at which $L(s_0)$ has a nontrivial kernel in the realizations of
\cref{sec:setting-realizations}. By \cref{lem:affine}, the roots of $L$ are the points $s_0$ with
$-s_0\in\sigma(L(0))$, and a Jordan chain of $L$ at $s_0$, that is $L(s_0)h_0=0$ and
$L(s_0)h_j=-h_{j-1}$, is a Jordan chain of the operator $L(0)$ at the eigenvalue $-s_0$. The
\emph{algebraic multiplicity} $\mult(s_0)$ is the dimension of the generalized eigenspace
$\bigcup_k\ker(L(0)+s_0)^k$. It is finite, and the roots are isolated, by \cref{prop:fredholm}.

\subsection{Floquet classes and the two sectors}\label{sec:setting-sectors}

A multifield $h$ is $4\pi$-periodic, so it contains all Fourier indices $m$; the background contains
only the indices allowed by~\eqref{eq:symmetries}. Accordingly, the space of profiles splits as
$X=X_A\oplus X_B$:
\begin{itemize}[leftmargin=*]
\item \emph{sector A}, with the symmetry of the background: $h_1,h_2,h_3$ with even $m$ and $h_4$ with
  odd $m$;
\item \emph{sector B}, with the opposite symmetry: $h_1,h_2,h_3$ with odd $m$ and $h_4$ with even $m$.
\end{itemize}
Since $\omegabar$ belongs to sector A, both sectors are invariant under $L(s)$ and $C(s)$; we write
$L_A,L_B$ and $C_A,C_B$ for the restrictions. In terms of the half period $\tau\mapsto\tau+2\pi$, sector A
consists of the profiles with the same symmetry as the background, and sector B of those with the opposite
one.

Multiplication by $e^{i\tau/2}$ is the shift $m\mapsto m+1$ of the Fourier index. It preserves the
spaces of \cref{sec:setting-realizations} up to a constant factor, commutes with every operation
in $\xi$ and with multiplication by functions of $\tau$, and conjugates $\partial_\tau$ to
$\partial_\tau+i/2$. Hence
\begin{equation}\label{eq:shift}
  e^{-i\tau/2}\,L(s)\,e^{i\tau/2}=L(s+\tfrac{i}{2}),\qquad
  e^{-i\tau/2}\,C(s)\,e^{i\tau/2}=C(s+\tfrac{i}{2}),
\end{equation}
and the shift exchanges the two sectors.

\begin{lemma}[Reduction to sector A]\label{lem:sectorB}
For every $s\in\C$,
\[
  \mult_L(s)=\mult_{L_A}(s)+\mult_{L_A}(s+\tfrac{i}{2}),
\]
and $L_A$ is $i$-periodic: $\mult_{L_A}(s+i)=\mult_{L_A}(s)$. The conjugation~\eqref{eq:shift} maps the
Jordan chains of $L_B$ at $s$ onto those of $L_A$ at $s+\tfrac i2$, and it maps the chains satisfying the
constraints onto chains satisfying the constraints.
\end{lemma}

\begin{proof}
The splitting $X=X_A\oplus X_B$ reduces $L(s)$, so multiplicities add. The shift by one Fourier index
maps $X_B$ onto $X_A$ and, by~\eqref{eq:shift}, conjugates $L_B(s)$ to $L_A(s+\tfrac i2)$ and $C_B(s)$
to $C_A(s+\tfrac i2)$; conjugation preserves kernels, Jordan chains and their lengths. The shift by two
indices maps $X_A$ to itself and gives the $i$-periodicity. The details for the realizations used below,
where the shift is an isometry up to the factor $\kappa_1$, are in \cref{app:lemmas}.
\end{proof}

Consequently the roots of $L$, counted per Floquet class $s+\tfrac{i}{2}\Z$, correspond to the roots of
$L_A$ counted modulo $i\Z$. All computations are done for $L_A$ in the strip
\[
  \widetilde Q=\{-\tfrac14<\Im s\le\tfrac34\},
\]
whose lower part $\{-\tfrac14<\Im s\le\tfrac14\}$ carries the sector-A roots (the \emph{A band}) and whose upper
half carries the sector-B roots, shifted by $\tfrac i2$ (the \emph{B band}). For instance, the root $\sB$
of the Main Theorem appears as $\sB+\tfrac i2$ in the B band.

\emph{Reality.} The coefficients of $\omegabar$ satisfy the reality condition
$v_{-m,-n}=\overline{v_{mn}}$~\cite{RT2019}; hence $\overline{L_A(\bar s)\,\bar h}=L_A(s)h$ for the
corresponding conjugation of coefficients, and the roots of $L_A$ are symmetric under $s\mapsto\bar s$.
For $L_B$ the same argument, combined with the shift, gives the symmetry $s\mapsto\bar s+i$ in the B band
(\cref{app:lemmas}). These symmetries are what make the four roots of the Main Theorem real.

\subsection{Constraint propagation}\label{sec:setting-constraints}

The equations $\Omega^+=0$ are not independent. Reiterer and Trubowitz show that if $F(\mu,\omega)=0$,
then $\omega^\sharp=F^\sharp(\mu,\omega)$ satisfies the linear homogeneous system
\begin{equation}\label{eq:sharp}
  O_\mu^\sharp\,\omega^\sharp+\mu\,R_\xi\Gamma_1^\sharp\,\omega^\sharp+\Gamma_2^\sharp(\omega,\omega^\sharp)=0,
\end{equation}
where $O_\mu^\sharp=\partial_\tau+\mu\operatorname{diag}(\xi\partial_\xi+1,(1+\xi)\partial_\xi+1)$, and
$\Gamma^\sharp_1,\Gamma^\sharp_2$ are again algebraic~\cite[\S3]{RT2019}. Linearizing this identity at
$(\mu,\omegabar)$, without assuming $F(\mu,\omega)=0$, gives the following lemma, proved in
\cref{app:lemmas}.

\begin{lemma}[Propagation of the constraints]\label{lem:KCML}
Let
\[
  K(s)=J^\sharp+B^\sharp+s\,I,\qquad J^\sharp=O^\sharp_\mu,\qquad
  B^\sharp=\mu\,R_\xi\Gamma_1^\sharp+\Gamma_2^\sharp(\omegabar,\cdot\,).
\]
There is an explicit affine family $M(s)=M_0+sM_1$ (\cref{app:formulas}) such that
\begin{equation}\label{eq:KCML}
  K(s)\,C(s)=M(s)\,L(s)
\end{equation}
as closed operators from the domain of $L$ in $Y(a,b)$ to the domain of $K$ in $Y(a',b')$, for all weights
$1<a'<a$, $1<b'<b$ (\cref{sec:setting-realizations}).
\end{lemma}

In particular, if $L(s)h=0$ and $K(s)$ is injective on its domain in the smaller space, then $C(s)h=0$:
an eigenvector of $L$ satisfies the constraints at every $s$ which is not a root of $K$. The loss of
radius from $(a,b)$ to $(a',b')$ is harmless, because all our certificates for $K$ are formulated in the
smaller space $(129/128,9/8)$.

\subsection{Realizations and the Fredholm property}\label{sec:setting-realizations}

We use two families of coefficient spaces. For weights $a,b>1$:
\begin{itemize}[leftmargin=*]
\item $Y(a,b)$ is the weighted $\ell^1$ space of~\cite{RT2019}, with norm
  $\sum_i\eta_i\sum_{m,n}a^{|m|}b^{|n|}|(h_i)_{mn}|$; the component weights $\eta_i>0$ give equivalent
  norms and only affect constants (\cref{app:lemmas}). We write $\Yp=Y(65/64,5/4)$ for the space of the
  existence proof.
\item $\mathbb T(a,b)$, the \emph{torus space}, is the weighted $\ell^2$ space with \emph{signed} Fourier
  index,
  \[
    \|h\|^2_{\mathbb T(a,b)}=\sum_i\sum_{m,n\in\Z}a^{2m}b^{2n}\,|(h_i)_{mn}|^2 .
  \]
  It is $L^2$ of the torus $\{|w|=a\}\times\{|z|=b\}$ in the variables $w=e^{i\tau/2}$, $z=e^{i\theta}$,
  $\xi=(z+z^{-1})/2$. On it, multiplication by an analytic function, the reflection $P$ and the division by
  $\xi$ act pointwise, which makes several bounds sharp (\cref{sec:counting}).
\end{itemize}
The inclusion $Y(a,b)\subset\mathbb T(a,b)$ has norm at most one.

\begin{proposition}[Fredholm property]\label{prop:fredholm}
In each of the spaces $Y(a,b)$ and $\mathbb T(a,b)$ with $1\le a\le65/64$ and $1<b\le5/4$:
\begin{enumerate}[label=\textup{(\roman*)},leftmargin=*]
\item $J$ is closed on its maximal domain $\Dom$, which coincides with the closure of the trigonometric
  polynomials in the graph norm, and $J+z$ is invertible for real $z\ge0$, with compact inverse;
\item $B$ is bounded, so $L(s)=J+B+s$ is closed on the fixed domain $\Dom$ for all $s$;
\item with $Q=(J+z_0)^{-1}$ for a fixed $z_0>0$, the family $L(s)Q=I+(B+s-z_0)Q$ is an analytic
  family of Fredholm operators of index zero, invertible for real $s$ large.
\end{enumerate}
Consequently the roots of $L$ are isolated and have finite algebraic multiplicity. The same holds for
$K(s)$ in the corresponding sharp spaces.
\end{proposition}

The proof is in \cref{app:lemmas}. The key point of~(i) is that the homogeneous solutions of $J+z$ in each
Fourier mode are multiples of $(1+\xi)^{-1-(z+im/2)/\mu}$ or $\xi^{-1-(z+im/2)/\mu}$, which are singular at
$\xi=-1$ or $\xi=0$, inside the Bernstein ellipse; so $J+z$ is injective on the maximal domain.

The certificates of \cref{sec:counting,sec:roots} are computed in the torus space, while the existence
proof, and with it the physical interpretation, lives in $\Yp$. The following lemma shows that the
choice does not matter in the region where we count.

\begin{lemma}[Independence of the realization]\label{lem:Lreal}
For $|s|\le 3.1$, the kernels and the Jordan chains of $L(s)$ in $\Yp$ and in $\mathbb T(65/64,5/4)$
coincide. Hence, in the disc $|s|\le3.1$, the roots of $L$ and their algebraic multiplicities are the same
in the two realizations.
\end{lemma}

\begin{proof}[Sketch]
Write $H(s)=I+(B+s)Q$ with $Q=J^{-1}$, the same operator in both spaces, and split the coefficients into a
finite box $G=\{|m|<1024,\ n<4096\}$ and its complement $T=I-G$. In both spaces
$\|T(H(s)-I)T\|<1$ for $|s|\le3.1$: the bound is $0.371$ in the torus, from a bound
$\|B\|\le3.9642$ certified through the symbol of $B$ (\cref{sec:counting}), and $0.2666$ in $\Yp$, from a bound on the
background part computed from the data and the published ball of~\cite{RT2019} (\cref{app:recovery}); both use closed forms for $\|QT\|$. If $H(s)x=0$ with $x$ in the
torus space, then $Gx$ is a finite vector, $Tx$ is the unique torus solution of
$THT\,Tx=-THGx$, and the Neumann series solves the same equation in $\Yp$; by uniqueness the two
solutions agree, so $x\in\Yp$. Chains are handled by induction, since the right-hand side
$-h_{j-1}$ is already in the strong space. The converse is the inclusion. Details, and the rational
verification of the constants, are in \cref{app:lemmas}.
\end{proof}

\begin{remark}\label{rem:radius}
The same \emph{radius recovery} argument shows more: an element of the kernel, or of a Jordan chain, in
a space of smaller radius automatically lies in the space of larger radius, for every $s$
(\cref{sec:physical}, \cref{lem:recovery}). It is used in \cref{sec:physical} to show that physical modes,
which are only assumed analytic in some unspecified neighbourhood, produce eigenvectors in $\Yp$.
\end{remark}
